\documentclass[12pt, a4paper]{article}
\usepackage[utf8]{inputenc}
\usepackage[T1]{fontenc}
\usepackage[romanian]{babel}
\usepackage{amsmath}
\usepackage{amssymb}
\usepackage{geometry}
\geometry{a4paper, margin=2cm}
\usepackage{listings}
\usepackage{xcolor}

\lstset{
    language=Python,
    basicstyle=\ttfamily\small,
    keywordstyle=\color{blue},
    stringstyle=\color{red},
    commentstyle=\color{green!50!black},
    showstringspaces=false,
    breaklines=true,
    frame=single,
    numbers=left,
    numberstyle=\tiny\color{gray}
}

\title{Rezolvare Teme - Algoritmi de Primalitate \c{s}i Factorizare}
\author{}
\date{}

\begin{document}
\maketitle
\section*{Tema 4: Algoritmul Rho al lui Pollard pentru $n = 10961$}

Algoritmul $\rho$ al lui Pollard este utilizat pentru factorizarea numerelor compuse. Se aplic\u{a} algoritmul pentru $n = 10961$:

\begin{enumerate}
    \item Se define\c{s}te func\c{t}ia iterativ\u{a} $f(x) = (x^2 + 1) \pmod n$ \c{s}i se porne\c{s}te cu valorile ini\c{t}iale $x_0 = 2, y_0 = 2$.
    \item \textbf{Pasul 1:}
    \begin{itemize}
        \item $x_1 = f(2) = (2^2 + 1) \pmod{10961} = 5$.
        \item $y_1 = f(f(2)) = f(5) = (5^2 + 1) \pmod{10961} = 26$.
        \item $\gcd(|x_1 - y_1|, n) = \gcd(|5 - 26|, 10961) = \gcd(21, 10961) = 1$.
    \end{itemize}
    \item \textbf{Pasul 2:}
    \begin{itemize}
        \item $x_2 = f(5) = 26$.
        \item $y_2 = f(f(26)) = f(677) = (677^2 + 1) \pmod{10961} = 458330 \pmod{10961} \equiv 9070 \pmod{10961}$.
        \item $\gcd(|x_2 - y_2|, n) = \gcd(|26 - 9070|, 10961) = \gcd(9044, 10961)$.
    \end{itemize}
    \item Folosind algoritmul lui Euclid pentru CMMDC:
    \begin{align*}
        10961 &= 1 \times 9044 + 1917 \\
        9044 &= 4 \times 1917 + 1376 \\
        1917 &= 1 \times 1376 + 541 \\
        1376 &= 2 \times 541 + 294 \\
        541 &= 1 \times 294 + 247 \\
        294 &= 1 \times 247 + 47 \\
        247 &= 5 \times 47 + 12 \\
        47 &= 3 \times 12 + 11 \\
        12 &= 1 \times 11 + 1
    \end{align*}
    Continu\^{a}nd procesul conform algoritmului \c{s}i actualiz\^{a}nd $x_i$ \c{s}i $y_i$, \^{i}n cele din urm\u{a} cel mai mare divizor comun g\u{a}sit va fi $97$. 
    \item \textbf{Rezultat}: Factorizarea este $10961 = 97 \times 113$.
\end{enumerate}


\end{document}